\documentclass[12pt,letterpaper]{article}
\usepackage{graphicx,textcomp}
\usepackage{natbib}
\usepackage{setspace}
\usepackage{fullpage}
\usepackage{color}
\usepackage[reqno]{amsmath}
\usepackage{amsthm}
\usepackage{fancyvrb}
\usepackage{amssymb,enumerate}
\usepackage[all]{xy}
\usepackage{endnotes}
\usepackage{lscape}
\newtheorem{com}{Comment}
\usepackage{float}
\usepackage{hyperref}
\newtheorem{lem} {Lemma}
\newtheorem{prop}{Proposition}
\newtheorem{thm}{Theorem}
\newtheorem{defn}{Definition}
\newtheorem{cor}{Corollary}
\newtheorem{obs}{Observation}
\usepackage[compact]{titlesec}
\usepackage{dcolumn}
\usepackage{tikz}
\usetikzlibrary{arrows}
\usepackage{multirow}
\usepackage{xcolor}
\newcolumntype{.}{D{.}{.}{-1}}
\newcolumntype{d}[1]{D{.}{.}{#1}}
\definecolor{light-gray}{gray}{0.65}
\usepackage{url}
\usepackage{listings}
\usepackage{color}
\usepackage{booktabs}
\usepackage{tabularx}
\usepackage[english]{babel}


\definecolor{codegreen}{rgb}{0,0.6,0}
\definecolor{codegray}{rgb}{0.5,0.5,0.5}
\definecolor{codepurple}{rgb}{0.58,0,0.82}
\definecolor{backcolour}{rgb}{0.95,0.95,0.92}

\lstdefinestyle{mystyle}{
	backgroundcolor=\color{backcolour},   
	commentstyle=\color{codegreen},
	keywordstyle=\color{magenta},
	numberstyle=\tiny\color{codegray},
	stringstyle=\color{codepurple},
	basicstyle=\footnotesize,
	breakatwhitespace=false,         
	breaklines=true,                 
	captionpos=b,                    
	keepspaces=true,                 
	numbers=left,                    
	numbersep=5pt,                  
	showspaces=false,                
	showstringspaces=false,
	showtabs=false,                  
	tabsize=2
}
\lstset{style=mystyle}
\newcommand{\Sref}[1]{Section~\ref{#1}}
\newtheorem{hyp}{Hypothesis}

\title{Problem Set 2 \\
	\large Applied Stats/Quant Methods }
\date{Due: October 23, 2025}
\author{Matilda Tomatis}


\begin{document}
	\maketitle
	\section*{Instructions}
\begin{itemize}
	\item Please show your work! You may lose points by simply writing in the answer. If the problem requires you to execute commands in \texttt{R}, please include the code you used to get your answers. Please also include the \texttt{.R} file that contains your code. If you are not sure if work needs to be shown for a particular problem, please ask.
	\item Your homework should be submitted electronically on GitHub.
	\item This problem set is due before 23:59 on Thursday October 23, 2025. No late assignments will be accepted.

\end{itemize}

	
	\vspace{.5cm}
	\section*{Question 1: Political Science}
		\vspace{.25cm}
	The following table was created using the data from a study run in a major Latin American city.\footnote{Fried, Lagunes, and Venkataramani (2010). ``Corruption and Inequality at the Crossroad: A Multimethod Study of Bribery and Discrimination in Latin America. \textit{Latin American Research Review}. 45 (1): 76-97.} As part of the experimental treatment in the study, one employee of the research team was chosen to make illegal left turns across traffic to draw the attention of the police officers on shift. Two employee drivers were upper class, two were lower class drivers, and the identity of the driver was randomly assigned per encounter. The researchers were interested in whether officers were more or less likely to solicit a bribe from drivers depending on their class (officers use phrases like, ``We can solve this the easy way'' to draw a bribe). The table below shows the resulting data.

\newpage
\begin{table}[h!]
	\centering
	\begin{tabular}{l | c c c }
		& Not Stopped & Bribe requested & Stopped/given warning \\
		\\[-1.8ex] 
		\hline \\[-1.8ex]
		Upper class & 14 & 6 & 7 \\
		Lower class & 7 & 7 & 1 \\
		\hline
	\end{tabular}
\end{table}

\begin{enumerate}
	
	\item [(a)]
	Calculate the $\chi^2$ test statistic by hand/manually (even better if you can do "by hand" in \texttt{R}).\\
	\vspace{1cm}
	
	\item [(b)]
	Now calculate the p-value from the test statistic you just created (in \texttt{R}).\footnote{Remember frequency should be $>$ 5 for all cells, but let's calculate the p-value here anyway.}  What do you conclude if $\alpha = 0.1$?\\
	
\vspace{1cm}
	\item [(c)] Calculate the standardized residuals for each cell and put them in the table below.
	\vspace{1cm}
	
	\begin{table}[h]
		\centering
		\begin{tabular}{l | c c c }
			& Not Stopped & Bribe requested & Stopped/given warning \\
			\\[-1.8ex] 
			\hline \\[-1.8ex]
			Upper class  &  &  &  \\
			\\
			Lower class &  &   &   \\
			
		\end{tabular}
	\end{table}
	
	
	\vspace{1cm}
	\item [(d)] How might the standardized residuals help you interpret the results?  
\end{enumerate}

\vspace{2cm}
\noindent\textbf{Answer a}\\
\vspace{0.5cm}
\noindent In order to compute the chi-squared statics we must compute the total sum for each row and for each column, as well as the grand total for the whole table. A quick way to do so is by "pivoting longer" the dataset, so it would still be tidy (i.e. there would still be a variable for each column, a requirement not respected let the datset be left wide and adding the row totals and the column totals to it) and from there computing the sum of observations grouping by "Class" and by "Obserevd outcome". 
\noindent After obtaining these necessary values, we can proceed to calculate the expected frequencies for each observation via the formula $\frac{\text{Row Total} \times \text{Column Total}}{\text{Grand Total}}$ . 
\vspace{0.5cm}
\lstinputlisting[language=R, firstline=79, lastline=80]{PS02_Tomatis.R}  
\vspace{0.5cm}
\noindent At this point every element necessary is obtained and the chi-squared value can be manually camputed. This is the reference formula for the computation: $\chi^2 = \sum \frac{(F_o - F_e)^2}{F_e}$
\vspace{0.5cm}
\lstinputlisting[language=R, firstline=84, lastline=84]{PS02_Tomatis.R} 
\begin{verbatim}
[1] "The Chi-squared value found is: 3.7912"
\end{verbatim}

\vspace{0.5cm}
\noindent\textbf{Answer b}\\
\vspace{0.5cm}
\noindent Once the test statistic is obtained we can calcualte the related p-value; to do so we must consider a chi-squared distribution with $(\text{number of rows} - 1) \times (\text{number of columns} - 1)$ degrees of freedom.
\vspace{0.5cm}
\lstinputlisting[language=R, firstline=91, lastline=92]{PS02_Tomatis.R}
\begin{verbatim}
	[1] "p-value: 0.1502 - 
	The variables don't showcase statistically significant association"
\end{verbatim}
\vspace{0.5cm}
\noindent Considering the p-value associated is bigger than the significance level alpha, we cannot reject our null hypothesis of statistical indipendence between the variables "Social Class of the driver" and "Observed bribery request". 

\vspace{0.5cm}
\noindent\textbf{Answer c}\\
\vspace{0.5cm}
\noindent We must now calcualte the standardized residuals for each cell of our data frame. The formula for the computation is the following: $z = \frac{F_o - F_e}{\sqrt{F_e \left(1 - \frac{R_{\text{tot}}}{G_T}\right)\left(1 - \frac{C_{\text{tot}}}{G_T}\right)}}$. 
\vspace{0.5cm}
\lstinputlisting[language=R, firstline=105, lastline=107]{PS02_Tomatis.R}
\noindent For the displaying of the results the tibble will be "pivoted wider" once again. 
\begin{table}[!htbp] \centering 
	\caption{Standardized Residuals for Class and Police Stop Outcomes} 
	\label{} 
	\begin{tabular}{@{\extracolsep{5pt}} ccccc} 
		\\[-1.8ex]\hline 
		\hline \\[-1.8ex] 
		& Class & Not stopped & Bribe requested & Stopped/given warning \\ 
		\hline \\[-1.8ex] 
		1 & Upper class & $0.322$ & $$-$1.642$ & $1.523$ \\ 
		2 & Lower class & $$-$0.322$ & $1.642$ & $$-$1.523$ \\ 
		\hline \\[-1.8ex] 
	\end{tabular} 
\end{table} 

\vspace{0.5cm}
\noindent\textbf{Answer d}\\
\vspace{0.5cm}
\noindent Standardized residulas are useful for the interpretation of a chi-squared coeffient, as they indicate theh direction in which the observed values diverge from the expected ones (computed under the assumption  of indipendence). The chi-squared coefficient alone does not show the magnitude of the distance of each observation, it might be inflated by a few larger deviations rather than by the whole sample. In our case, we can observe how none of the z-values exceed +/- 2, a general rule of thumb threshold beyond which it's plausible to question the indipendence of the variables under examination. 


\newpage

\section*{Question 2: Economics}
Chattopadhyay and Duflo were interested in whether women promote different policies than men.\footnote{Chattopadhyay and Duflo. (2004). ``Women as Policy Makers: Evidence from a Randomized Policy Experiment in India. \textit{Econometrica}. 72 (5), 1409-1443.} Answering this question with observational data is pretty difficult due to potential confounding problems (e.g. the districts that choose female politicians are likely to systematically differ in other aspects too). Hence, they exploit a randomized policy experiment in India, where since the mid-1990s, $\frac{1}{3}$ of village council heads have been randomly reserved for women. A subset of the data from West Bengal can be found at the following link: \url{https://raw.githubusercontent.com/kosukeimai/qss/master/PREDICTION/women.csv}\\

\noindent Each observation in the data set represents a village and there are two villages associated with one GP (i.e. a level of government is called "GP"). Figure~\ref{fig:women_desc} below shows the names and descriptions of the variables in the dataset. The authors hypothesize that female politicians are more likely to support policies female voters want. Researchers found that more women complain about the quality of drinking water than men. You need to estimate the effect of the reservation policy on the number of new or repaired drinking water facilities in the villages.
\vspace{.5cm}

\begin{table}[h!]
	\centering
	\caption{\footnotesize Names and description of variables from Chattopadhyay and Duflo (2004).}
	\label{tab:women_desc}
	\begin{tabularx}{\textwidth}{lX}
		\toprule
		\textbf{Name} & \textbf{Description} \\
		\midrule
		GP & An identifier for the Gram Panchayat (GP). \\
		village & Identifier for each village. \\
		reserved & Binary variable indicating whether the GP was reserved for women leaders or not. \\
		female & Binary variable indicating whether the GP had a female leader or not. \\
		irrigation & Number of new or repaired irrigation facilities in the village since the reserve policy started. \\
		water & Number of new or repaired drinking-water facilities in the village since the reserve policy started. \\
		\bottomrule
	\end{tabularx}
\end{table}



\newpage
\begin{enumerate}
	\item [(a)] State a null and alternative (two-tailed) hypothesis. 
	
	\vspace{1cm}
	\item [(b)] Run a bivariate regression to test this hypothesis in \texttt{R} (include your code!).
	
	\vspace{1cm}
	\item [(c)] Interpret the coefficient estimate for reservation policy. 
\end{enumerate}

\vspace{2cm}
\noindent\textbf{Answer a}\\
\vspace{0.5cm}
\begin{enumerate}
	\item $H_0$ states: $\beta_1 = 0$, interpretable as having reserved positions for female leaders has no impact on the creation and repair of irrigation facilities. \\
	\item $H_a$ states: $\beta_1 \neq 0$, hence part of the leading positions being reserved for women does have an impact on the creation and repair of irrigation facilities. \\
\end{enumerate}

\vspace{0.5cm}
\noindent\textbf{Answer b}\\
\vspace{0.5cm}
\noindent We will proceed to test our hyphothesis by running the following bivariate regression: 
\noindent $\text{Irrigation} = \beta_0 + \beta_1 \text{Reserved} + \varepsilon$
\lstinputlisting[language=R, firstline=136, lastline=136]{PS02_Tomatis.R}
\begin{table}[!htbp] \centering 
	\caption{Bivariate regression model with GP as Y and irrigation as X1} 
	\label{} 
	\begin{tabular}{@{\extracolsep{5pt}}lc} 
		\\[-1.8ex]\hline 
		\hline \\[-1.8ex] 
		& \multicolumn{1}{c}{\textit{Dependent variable:}} \\ 
		\cline{2-2} 
		\\[-1.8ex] & irrigation \\ 
		\hline \\[-1.8ex] 
		reserved & $-$0.369 \\ 
		& (1.122) \\ 
		& \\ 
		Constant & 3.388$^{***}$ \\ 
		& (0.650) \\ 
		& \\ 
		\hline \\[-1.8ex] 
		Observations & 322 \\ 
		R$^{2}$ & 0.0003 \\ 
		Adjusted R$^{2}$ & $-$0.003 \\ 
		Residual Std. Error & 9.506 (df = 320) \\ 
		F Statistic & 0.108 (df = 1; 320) \\ 
		\hline 
		\hline \\[-1.8ex] 
		\textit{Note:}  & \multicolumn{1}{r}{$^{*}$p$<$0.1; $^{**}$p$<$0.05; $^{***}$p$<$0.01} \\ 
	\end{tabular} 
\end{table} 
\vspace{0.5cm}
\noindent\textbf{Answer c}\\
\vspace{0.5cm}
\noindent The $\beta_1$ estimate has a value of -0.369, indicating that when a Gram Panchayat (GP) has positions reserved for female leaders the number of new or repaired irrigation facilities in its villages decreases by approximately one-third, on average, compared to non-reserved GPs. The standard error estimate is 1.122.  

\noindent Across all conventional significance thresholds (0.10, 0.05 and 0.01), the coefficient computed shows no statistically significant impact on the creation and restoration of irrigation facilities. However, the model might suffer from oversimplification: it is difficult to hypothesize that only the reservation policy might have an impact on irrigation facility funding, and our estimate might suffer from omitted variable bias.
%\end{enumerate}
\end{document}
